\section{Method}
\label{sec:method}
\subsection{Setting}
\label{sec:setting}
\textbf{Model and agent.} The model is Gemma 4 E4B~\citep{Team2026Gemma4}: 42 layers, residual width $n = 2560$, MLP width 10240, activation \texttt{gelu\_pytorch\_tanh}. We use its INT4 quantisation-aware checkpoint, served with vLLM. The agent calls six tools: run a command, read, edit and write a file, status, and submit. It uses greedy decoding, a thinking budget of 2,560 tokens, at most 34 tool calls and 7.5 minutes per episode. A task is solved when hidden tests pass. We edit the released base model (the \emph{unedited model}) and, in three experiments, our wave-5 RL adapter.

\textbf{Tasks and what is in-sample.} The 170 \emph{development tasks} are 106 competition-format tasks and 64 SWE-rebench~\citep{Badertdinov2025Swerebench} tasks. The SWE-rebench repositories are disjoint from all RL training pools. The 130 \emph{held-out tasks} come from 78 repositories. These share no repository with the development tasks and never appeared in the logs that we computed directions from (\S\ref{sec:heldout}). The diagnosis (\S\ref{sec:diagnosis}) used episodes of the 170 development tasks. We split logged decisions into training and held-out sets by episode, not by task. Development gains are therefore in-sample at the task level. The 130 held-out tasks are the out-of-sample test of solve rate.

\textbf{What each data split is used for.} The training split of logged decisions from development-task episodes fits directions and read-outs. The held-out split of those decisions (by episode, same tasks) checks direction AUC, sizes gated edits, measures realised shifts, and supports decision replay. Replay on this split therefore measures in-distribution control. The 170 development tasks serve diagnosis, agent evaluation, and the choice of strengths and of edits to repeat. The 130 held-out tasks serve only the two pre-registered tests.

\textbf{Unit of comparison.} An \emph{arm} is one model variant (the unedited model, an edit or a control) evaluated on every task. Greedy runs of this stack are not repeatable. A \emph{batch} is one run of every arm on all tasks at once. vLLM loads the base checkpoint and serves every arm as a LoRA module; the unedited model is an all-zero adapter. Each arm has its own evaluator process. The arms share one continuous-batching queue, so a request's batch neighbours vary between runs. Two copies of the unedited model in the same held-out run differed by 11.7 tasks-equivalent (\S\ref{sec:heldout}). We therefore make every comparison within one batch, against the unedited arm served alongside the edit. Every batch includes the unedited model. The primary unit of analysis is the batch. We report the mean per-batch difference in solved tasks with a 95\% $t$ confidence interval (CI) and a two-sided paired $t$ test over batches. Task-level tests (per-task differences summed over batches) are secondary, and we label them as such.

\subsection{Locating a direction}
\label{sec:locate}
A binary label on logged decision points defines a direction. For ``explore more before editing'', the label is whether the next tool call makes the episode's first edit (class 1) or explores further (class 0). The position is the first generated token of that call. We record the layer outputs at these positions for 400 training and 160 held-out decisions per class. For layer $l$, the direction $d_l = \mu_1 - \mu_0$ is the difference of the class means of the layer output on the training split. Its norm $\lVert d_l \rVert$ is the \emph{class gap}. We use this single split, without further regularisation or cross-validation. We keep a direction if projecting held-out decisions onto $d_l$ separates the classes, measured as AUC. For explore-more, the AUC is 0.77, 0.81 and 0.83 at layers 18, 20 and 22, the three layers we write. To make the agent explore more, we write $-d_l$.

\subsection{The primitive: a hand-written neuron}
\label{sec:primitive}
\textbf{Why weights and not a hook.} The competition scores a submitted checkpoint served by a stock inference engine. An intervention must therefore live in the weights; an activation hook cannot be submitted. Gemma 4 has no bias terms, so we must also build a constant shift from the input. In the MLP of layer $l$, with $\hat{x} \in \mathbb{R}^n$ the normalised MLP input, neuron $j$ computes
\begin{equation}
\label{eq:neuron}
    a_j = \gelu(\mathrm{gate}_j \cdot \hat{x}) \, (\mathrm{up}_j \cdot \hat{x}), \qquad y = \sum_j \mathrm{down}_{:,j} \, a_j, \qquad h_{\mathrm{out}} = s_l \left( x + w_{\mathrm{post}} \odot y / \mathrm{rms}(y) \right),
\end{equation}
where $\mathrm{gate}_j$ and $\mathrm{up}_j$ are rows of the two input projections, $\mathrm{down}_{:,j}$ is a column of the output projection, $w_{\mathrm{post}} \in \mathbb{R}^n$ is the weight of the output normalisation, $s_l$ is the layer's output scalar, $x$ is the residual input and $\odot$ is the elementwise product. The MLP output is RMS-normalised, so only its direction matters.

\textbf{Read.} We fit a ridge read-out $q$ with target 1, so that $q \cdot \hat{x} \approx 1$ at every position. The fit uses relative penalty $10^{-3}$, the dual form, and the 3,998 training-split positions (about five per training decision call). On held-out positions, $q \cdot \hat{x}$ has mean 0.98--0.99 and standard deviation (SD) 0.02--0.04 across the three layers for the always-on edit (other builders: Appendix~\ref{app:gated}).

\textbf{Neuron choice.} In each layer we overwrite the neuron with the smallest product of mean absolute activation and down-column norm (0.22--0.32 of the layer median). Overwriting it removes almost nothing.

\textbf{Write.} The overwritten neuron gets
\begin{equation}
\label{eq:write}
    \mathrm{gate}_j = G\, q, \qquad \mathrm{up}_j = q, \qquad \mathrm{down}_{:,j} = \pm\,\beta \, \rho \, \sqrt{n} \; \hat{u} \, / \, \gelu(G), \qquad \hat{u} = \unit(d_l / w_{\mathrm{post}}),
\end{equation}
with gain $G = 4$, $\rho$ the mean RMS of the natural MLP output $y$ on training positions, $\unit(v) = v / \lVert v \rVert$, the division elementwise, and the sign chosen by the author. The neuron fires at $a_j \approx \gelu(G)$ everywhere. It adds a vector of norm $\beta \rho \sqrt{n}$ along $\hat{u}$, which is $\beta$ times the norm of the natural output. The map $d_l \mapsto \hat{u}$ undoes the output normalisation: $w_{\mathrm{post}} \odot \hat{u} = \eta\, d_l / \lVert d_l \rVert$ with $\eta = \lVert w_{\mathrm{post}} \odot \hat{u} \rVert$.

\textbf{Predicted shift and the strength scale.} Assume the natural output $y$ is orthogonal to $\hat{u}$ and has RMS $\rho$. Then $y + \beta \rho \sqrt{n}\, \hat{u}$ has RMS $\rho \sqrt{1+\beta^2}$, and the full change of the layer output is
$\Delta h = s_l \big[ \tfrac{\beta \sqrt{n}}{\sqrt{1+\beta^2}}\, w_{\mathrm{post}} \odot \hat{u} - \big(1 - \tfrac{1}{\sqrt{1+\beta^2}}\big)\, w_{\mathrm{post}} \odot y / \rho \big]$.
The second term shrinks the natural write by $1 - 1/\sqrt{1+\beta^2} \approx \beta^2/2$. This is 0.04--1.2\% for the $\beta$ of strengths up to 0.25 and 14.3\% at the cap $\beta = 0.6$ (Appendix~\ref{app:params}). Eq.~(\ref{eq:shift}) keeps only the first term. It ignores the part of the natural write that already lies along $d_l$; the shrinkage acts on that part too, so at layer 22 at the cap the natural write's component along $d_l$ also shrinks by 14\%. Also, because $\hat{u} \propto d_l / w_{\mathrm{post}}$, the pre-normalisation inner product $y \cdot \hat{u}$ weights each coordinate of $(w_{\mathrm{post}} \odot y) \cdot d_l$ by $1/w_{\mathrm{post}}^2$. Orthogonality of $y$ and $\hat{u}$ therefore makes $w_{\mathrm{post}} \odot y$ and $w_{\mathrm{post}} \odot \hat{u}$ orthogonal only when $w_{\mathrm{post}}$ is constant. Under these approximations, the layer output moves along $d_l / \lVert d_l \rVert$ by
\begin{equation}
\label{eq:shift}
    \Delta_l(\beta) = s_l \cdot \frac{\beta}{\sqrt{1+\beta^2}} \cdot \sqrt{n} \cdot \eta .
\end{equation}
This predicted shift is the projection of the change in $h_{\mathrm{out}}$ onto the unit direction. We report it in units of the class gap. An edit of \emph{strength} $t$ sets $\beta$ per layer so that $\Delta_l(\beta) = t\, \lVert d_l \rVert$. With $\kappa = t \lVert d_l \rVert / (s_l \sqrt{n}\, \eta)$, this gives $\beta = \kappa / \sqrt{1 - \kappa^2}$, capped at 0.6. Strength 0.25 gives $\beta = 0.057$, 0.060 and 0.154 at the three layers (Appendix~\ref{app:params}). The layer scalar enters through $\kappa$. At layer 22 ($s_{22} = 0.173$), $\kappa = 0.61\,t$. The shift therefore saturates at $t \approx 1.64$ ($\kappa = 1$), and the cap binds from $t \approx 0.84$ (layers 18 and 20: from $t \approx 2.1$--2.3). We measure the realised shift on held-out positions wherever a build report records it. The prediction holds closely for always-on edits. Over 12 edit--layer pairs built by Claude and Gemini, the realised shift is 0.96--1.14 times Eq.~(\ref{eq:shift}) (mean 1.06). For gated edits, the ratio is the realised shift averaged over held-out gate-ON positions, divided by Eq.~(\ref{eq:shift}) at the builder's recorded $\beta$. This ratio follows how strongly the gate fires there: 0.92--1.17 per unit of firing over 13 GPT and Gemini pairs (\texttt{scripts/shift\_check.py}). The raw ratio ranges from 0.64--0.71 for Claude's early-calls gate to 2.14 for a Gemini guard that fires at twice its reference level. For Claude's gates, the realised value is the bisection target itself. For the early-calls gate, this target is 0.3 gaps, averaged over all its held-out ON positions, where its mean activation is 2.0--2.2 of the 4.0 of full firing. The decision token of calls up to 4 moves more (Appendix~\ref{app:gated}).

\textbf{Gated neuron.} A gated edit fires only where a read direction for an agent state is present. Its gate row combines a read-out $r$ for the gating label with the constant read-out, $\mathrm{gate}_j \propto r - \theta q$. The pre-activation is therefore about $+G$ where the state is present and negative elsewhere. At negative pre-activations $\gelu$ is near zero but not zero. \texttt{gelu\_pytorch\_tanh} is negative on $(-3, 0)$, with minimum $-0.17$ at $-0.75$. A weakly negative pre-activation therefore writes up to 4\% of a full firing with reversed sign. \S\ref{sec:control} measures how often the gate fires outside its fit. The gated edits of the benchmark fit $r$ by ridge regression. They size the write by measurement on the held-out split instead of Eq.~(\ref{eq:shift}), so their sizing is in-sample for that split. Appendix~\ref{app:gated} gives these details and the other authors' rotation and attention edits. A neuron edit in three layers changes nine tensors: one gate row, one up row and one down column per layer. We store it as a LoRA update, exact up to bf16 rounding.

\textbf{Matched random control.} The control overwrites the same neurons with the same gate row and up row. We build its write direction in residual space. We take $d_l$ with a random sign per dimension and make it orthogonal to $d_l$ by Gram--Schmidt. We then map it through $1 / w_{\mathrm{post}}$ and normalise it to unit length, so the down column has the same norm as the edit's. Its cosine with $d_l$ is zero, and its size in the residual stream is 0.99--1.04 of the edit's per layer. This strict construction applies to the unconditional edits of \S\ref{sec:control}--\ref{sec:devgain}. In the three-author benchmark, each author built its own control (Appendix~\ref{app:gated}).

\textbf{The compiled edit closely reproduces the intended intervention.} We compare the adapter as served with an inference-time hook that adds the same neuron write to the MLP output. Both run teacher-forced through a twin of the serving numerics (INT4 weights, bfloat16) on 76 held-out decision calls. We test the always-on explore-more edit and the before-first-edit gate. Relative to the unedited model, the residual-stream changes of adapter and hook have cosine 0.972--0.992 and 0.984--0.996 across eight layers from 18 to 41. Their norm ratios are within 0.4\% of 1. Over completion tokens, KL(hook$\,\Vert\,$adapter) is $1.1\times10^{-4}$ and $1.2\times10^{-4}$. This is 99 and 110 times smaller than the edit's own effect, KL(unedited$\,\Vert\,$adapter) $= 0.011$ and $0.013$. At the decision token itself, KL(unedited$\,\Vert\,$adapter) is $3.4\times10^{-5}$ for both edits (Appendix~\ref{app:hook}). Most of the measured effect is therefore on completion tokens after the decision. A plain vector addition matches less closely (cosine 0.80 and 0.73 at layer 22; Appendix~\ref{app:hook}).

\subsection{The loop and the written-prediction protocol}
\label{sec:loop}
Figure~\ref{fig:loop} shows the loop. Algorithm~\ref{alg:loop} (Appendix~\ref{app:params}) lists its six steps: diagnose, locate, write, predict, verify by replay and by agent evaluation, and score. The author is an LLM with shell access to the model weights, stored activations and logged episodes. Steps 1--4 run on a CPU; only verification uses a GPU.

\textbf{Prediction protocol.} Step 4 produces a file that is fixed before step 5 starts. In the three-author benchmark (\S\ref{sec:arena}), each edit's entry has the fields \texttt{diagnosis}, \texttt{mechanism}, \texttt{predicted\_behaviour} and \texttt{predicted\_solved\_delta}. The last field holds a point value and an 80\% interval in solved tasks out of 170, relative to the unedited model of the same batch. A prediction is a hit when the first measured difference falls inside the interval.
